\section{Arithmetic}
\label{sec:arithmetic}

Arithmetic is the floor everything else stands on, and Mathilda takes it seriously.
Rather than one kind of number, it keeps a small \emph{tower} of numeric
representations and moves between them automatically:

\begin{itemize}
  \item \textbf{Machine integers} — signed 64-bit words, the hardware's native
        integer, used while a value is small enough to fit.
  \item \textbf{Arbitrary-precision integers} — GMP~\cite{gmp} ``bignums'', into
        which a machine integer is promoted the instant it would overflow, and out
        of which a result is demoted the instant it fits again.
  \item \textbf{Exact rationals} — pairs of exact integers, kept in lowest terms.
  \item \textbf{Machine reals} — IEEE-754 double-precision floating point, the fast
        approximate numbers of the hardware.
  \item \textbf{Arbitrary-precision reals} — MPFR~\cite{mpfr} values, correctly
        rounded to as many digits as you ask for.
\end{itemize}

The whole of this section is about how these five representations behave and how
they hand off to one another. We begin with the way numbers are typed and tested, and
the operators that combine them; then we take the representations in three groups: the
machine numbers (fast, fixed-size), the arbitrary-precision numbers (exact integers and
rationals, and MPFR reals), and finally interval arithmetic, which computes with
\emph{guaranteed enclosures} rather than points.

\subsection{The types of number}
\label{ssec:number-types}

In Mathilda every value carries a \emph{head} that names its type, and \B{Head} reads
it off. The four numeric heads are \texttt{Integer}, \texttt{Real}, \texttt{Rational},
and \texttt{Complex}:

\begin{codepairs}
\pair{arithmetic/heads/1}{A whole number has head \texttt{Integer}.}
\pair{arithmetic/heads/2}{So does a bignum---an arbitrary-precision integer is still an
\texttt{Integer}, not a separate type.}
\pair{arithmetic/heads/3}{A machine real has head \texttt{Real}.}
\pair{arithmetic/heads/4}{A fraction has head \texttt{Rational}: internally it is
\texttt{Rational[p, q]}.}
\pair{arithmetic/heads/5}{And a complex number has head \texttt{Complex}, stored as
\texttt{Complex[re, im]}.}
\end{codepairs}

Rather than test the head by hand, you use the \emph{predicates}---the functions ending
in \texttt{Q} (for ``question'') that return \texttt{True} or \texttt{False}. Mathilda
provides a full family for numbers:

\begin{center}\small
\begin{tabular}{@{}ll@{}}
\toprule
Predicate & True when the argument is\ldots \\
\midrule
\B{IntegerQ}        & an integer (\texttt{Integer} or a bignum) \\
\B{MachineIntegerQ} & a machine-word (64-bit) integer, but not a bignum \\
\B{RationalQ}       & an exact rational (an integer or a \texttt{Rational}) \\
\B{ComplexQ}        & a \texttt{Complex} number \\
\B{ExactNumberQ}    & an exact number (integer, rational, or exact complex) \\
\B{InexactNumberQ}  & an inexact number (a machine or MPFR real) \\
\B{NumberQ}         & any explicit number \\
\B{NumericQ}        & a number \emph{or} a numeric constant such as $\pi$ \\
\bottomrule
\end{tabular}
\end{center}

\begin{codepairs}
\pair{arithmetic/predicates/1}{\B{IntegerQ} accepts any integer.}
\pair{arithmetic/predicates/2}{\B{MachineIntegerQ} is stricter: a bignum, having
overflowed the machine word, is not a machine integer.}
\pair{arithmetic/predicates/3}{\B{RationalQ} recognises a fraction.}
\pair{arithmetic/predicates/4}{A Gaussian integer is an \emph{exact} number, so
\B{ExactNumberQ} is true.}
\pair{arithmetic/predicates/5}{A fifty-digit MPFR value is \emph{inexact}, so
\B{InexactNumberQ} is true.}
\pair{arithmetic/predicates/6}{\B{ComplexQ} tests for the \texttt{Complex} head.}
\end{codepairs}

Each predicate is a small structural test: it looks at the head and type of its one
argument and returns \texttt{True} or \texttt{False}, never leaving the call
unevaluated for a concrete number. Every one of these functions carries its own
documentation, which you can read at the prompt with \texttt{?Name}; here, for
instance, is what \B{MachineIntegerQ} reports about itself:

\usagebox{MachineIntegerQ}

Throughout the rest of this section, boxes like the one above reproduce a function's
own docstring---exactly the text \texttt{?Name} prints---together with its attributes.

\subsection{The arithmetic operators}
\label{ssec:operators}

The operators you type---$+$, $-$, $\times$, $/$, $\wedge$---are surface syntax for a
few built-in heads, and \emph{only a few}. \B{Plus} and \B{Times} and \B{Power} are the
primitives; subtraction, division, and negation are not separate operations but
rewrites into them, as \B{FullForm} reveals:

\begin{codepairs}
\pair{arithmetic/operator-forms/1}{Subtraction is \B{Plus} against a negated term:
$a-b$ is \texttt{Plus[a, Times[-1, b]]}.}
\pair{arithmetic/operator-forms/2}{Division is \B{Times} against a reciprocal power:
$a/b$ is \texttt{Times[a, Power[b, -1]]}.}
\pair{arithmetic/operator-forms/3}{Unary minus is multiplication by $-1$.}
\pair{arithmetic/operator-forms/4}{Even $\sqrt{x}$ is a \B{Power}, with exponent the
exact rational $1/2$.}
\end{codepairs}

Reducing everything to \B{Plus}, \B{Times}, and \B{Power} is what keeps the evaluator
small: a single canonicalizer for each, driven by attributes\calloutnote{Under the
hood}{The attributes do the organising. \texttt{Flat} means nested calls are spliced
(\texttt{Plus[a, Plus[b, c]]} becomes \texttt{Plus[a, b, c]}); \texttt{Orderless} sorts
the arguments into a canonical order so that \texttt{a + b} and \texttt{b + a} are the
same expression and like terms sit together to be combined; \texttt{OneIdentity} lets a
one-argument call collapse; \texttt{Listable} threads the operation over lists; and
\texttt{NumericFunction} marks it for numeric evaluation. Underneath, the exact-integer
paths add and multiply in a $128$-bit intermediate and, on overflow, fall through to
GMP~\cite{gmp}, promoting the result to a bignum---so the answer is always exact, and
the fast machine-word path is taken whenever it fits.}, handles addition,
subtraction, multiplication, division, powers, and roots at once. Here are the three,
with their own docstrings and attributes:

\usagebox{Plus}
\usagebox{Times}
\usagebox{Power}

\subsection{Machine-precision computation}
\label{ssec:machine-precision}

At the bottom of the tower are the machine numbers: a signed 64-bit integer and an
IEEE-754 double. These are the types the processor implements directly, and Mathilda
uses them whenever it can, because they are fast.

An integer stays a machine word only while it fits in 64 bits. The largest such value
is $2^{63}-1$; one step beyond and Mathilda promotes it to an arbitrary-precision
integer, silently and without loss:

\begin{codepairs}
\pair{arithmetic/machine-integers/1}{$2^{62}$ fits in a signed 64-bit word and stays
a machine integer.}
\pair{arithmetic/machine-integers/2}{$2^{63}$ is one past the largest machine integer,
so it is promoted to a GMP bignum---still exact.}
\pair{arithmetic/machine-integers/3}{Addition promotes too: two machine integers whose
sum overflows produce a bignum. The promotion is detected with a wider (128-bit)
intermediate.}
\end{codepairs}

Machine reals are IEEE-754 binary64 (``double'') values. Their characteristics are
reported by a handful of read-only system constants, derived at start-up from the
platform's \texttt{<float.h>}:

\begin{codepairs}
\pair{arithmetic/machine-limits/1}{A double carries about $15.95$ decimal digits (a
53-bit mantissa).}
\pair{arithmetic/machine-limits/2}{\B{\$MachineEpsilon}: the gap between $1$ and the
next representable double.}
\pair{arithmetic/machine-limits/3}{The largest finite double.}
\pair{arithmetic/machine-limits/4}{And the smallest positive normal double.}
\end{codepairs}

Here is the single most important thing to understand about Mathilda's machine reals,
and a real difference from Mathematica: \textbf{Mathilda does not use significance
arithmetic}.\calloutnote{Theory}{\emph{Significance arithmetic} (which Mathematica uses
for arbitrary-precision numbers) attaches to each number an estimate of its error and
propagates that estimate through every operation, so that catastrophic cancellation
visibly reduces the reported precision. It is elegant but only a first-order model of
error, and it makes a number's behaviour depend on its history. Mathilda takes the other
well-trodden road: a machine real is a fixed-precision IEEE double, exactly as specified
by the floating-point standard~\cite{ieee754}. If you need to know how error propagates,
you either reason about it yourself---the classic reference is Goldberg~\cite{goldberg1991}---or
reach for the rigorous tool in \S\ref{ssec:interval}, interval arithmetic.} A machine
real is just an IEEE double, with a fixed 53-bit mantissa; it
carries no error term, and operations do not track how many digits remain trustworthy.
What you get is exactly what the hardware, C, or NumPy would give you---no more and no
less:

\begin{codepairs}
\pair{arithmetic/machine-doubles/1}{Dividing two machine reals gives a machine real,
printed to six figures by default.}
\pair{arithmetic/machine-doubles/2}{$0.1+0.2$ \emph{prints} as \texttt{0.3}, because
six figures round to that.}
\pair{arithmetic/machine-doubles/3}{But it is not exactly $0.3$: the residual is the
true IEEE rounding error, shown plainly rather than hidden.}
\pair{arithmetic/machine-doubles/4}{The precision of any machine real is the symbol
\texttt{MachinePrecision}.}
\end{codepairs}

This fixed-precision model is not a limitation so much as a deliberate choice, and for
the overwhelming majority of numerical work---linear algebra, differential equations,
optimization, plotting---it is exactly right: it is fast, it is predictable, and it
matches every other serious numerical system bit for bit. When you genuinely need more
than a double can hold, you move up the tower.

\subsection{Arbitrary-precision computation}
\label{ssec:arbitrary-precision}

Exactness in Mathilda is never an option you switch on; it is the default, and it never
silently degrades. Integers grow without bound, factorials and powers stay exact, and
the arithmetic behind them is GMP~\cite{gmp}, the same library that underpins much of
computational number theory:

\begin{codepairs}
\pair{arithmetic/bignums/1}{A thirty-one-digit power, kept exactly.}
\pair{arithmetic/bignums/2}{$100!$ is a $158$-digit integer, computed exactly by GMP.}
\pair{arithmetic/bignums/3}{Factoring the Mersenne number $2^{67}-1$---the number Cole
famously factored by hand in 1903.}
\end{codepairs}

Rational numbers are exact too. A ratio of integers is kept as a reduced fraction, its
numerator and denominator sharing no common factor, and it collapses to an integer the
moment its denominator becomes~$1$:

\begin{codepairs}
\pair{arithmetic/rationals/1}{A sum of fractions is another fraction, reduced to lowest
terms.}
\pair{arithmetic/rationals/2}{Reduction happens on construction, so $2/4$ is already
$1/2$.}
\pair{arithmetic/rationals/3}{When the denominator reduces to $1$, the rational
collapses to a plain integer.}
\end{codepairs}

For \emph{approximate} numbers beyond machine precision, Mathilda uses MPFR, which
computes correctly-rounded floating point to any requested precision. You reach it with
the two-argument form \B{N}\texttt{[expr, d]}, asking for \texttt{d} digits:

\begin{codepairs}
\pair{arithmetic/mpfr-N/1}{$\pi$ to fifty digits, via MPFR's \texttt{mpfr\_const\_pi}.}
\pair{arithmetic/mpfr-N/2}{$e$ to a hundred digits.}
\pair{arithmetic/mpfr-N/3}{$\sqrt2$ to forty digits---an MPFR power, not a table
lookup.}
\pair{arithmetic/mpfr-N/4}{Ap\'ery's constant $\zeta(3)$ to fifty digits. The zeta
function is evaluated by FLINT's rigorous ball arithmetic~\cite{flint,arb}, not by
MPFR.}
\end{codepairs}

So the numeric substrate is really three libraries working together: GMP for exact
integers, MPFR for correctly-rounded arbitrary-precision reals, and FLINT/Arb for the
harder special functions, where rigorous complex-ball arithmetic guarantees every
printed digit. \B{N} orchestrates them, replacing each exact leaf of an expression with
a numeric value of the requested precision and letting the ordinary evaluator re-run.

A number's precision is itself a value you can query, and here Mathilda is honest about
a small subtlety:

\begin{codepairs}
\pair{arithmetic/precision/1}{Exact numbers have \emph{infinite} precision.}
\pair{arithmetic/precision/2}{But \B{N}\texttt{[Pi, 50]} reports about $50.27$ digits,
not $50$: fifty decimal digits are stored as $\lceil 50\log_2 10\rceil = 167$ bits,
which is a hair more than fifty digits' worth.}
\pair{arithmetic/precision/3}{You can also state a number's precision on input: the
literal \texttt{1.5`20} is $1.5$ carried at twenty digits.}
\end{codepairs}

Precision is \emph{contagious}, and when representations of different precision meet, the
lowest wins---machine precision being the floor. This, too, follows Mathematica:

\begin{codepairs}
\pair{arithmetic/contagion/1}{Add a machine real to a hundred-digit number and the
result drops to machine precision: one approximate double poisons the lot.}
\pair{arithmetic/contagion/2}{Applying \B{N} to an already-precise number preserves its
precision rather than collapsing it to a double.}
\pair{arithmetic/contagion/3}{$1001!$ has $2{,}571$ digits, far past a double's range,
so \B{N} keeps it as a machine-precision value with an unbounded exponent---a real with
a mantissa but no overflow.}
\end{codepairs}

\subsection{Interval arithmetic}
\label{ssec:interval}

The tools so far give you either an exact answer or an approximate one whose error you
must reason about yourself. \emph{Interval arithmetic} gives you a third thing: an
approximate answer that comes with a \emph{proof}. Instead of a single number, you
compute with an interval $[a,b]$ that is guaranteed to contain the true value, and every
operation returns an interval guaranteed to contain the true result. The idea goes back
to Moore~\cite{moore2009}; Mathilda implements it as a first-class numeric object,
\B{Interval}\texttt{[\{lo, hi\}]}.\calloutnote{Pitfall}{Mathilda's intervals differ from
Mathematica's in two deliberate ways, both documented. A machine number becomes a
\emph{point} interval---\texttt{Interval[1.]} is \texttt{Interval[\{1., 1.\}]}---rather
than being preloaded with a half-ULP uncertainty band; if you want an uncertainty
propagated, you supply it. And \B{Interval} is a set, not a scalar, so it is deliberately
neither \texttt{Listable} nor a \texttt{NumericFunction}. The enclosures are rigorous in
the spirit of the interval standard, though Mathilda implements no formal decoration
system.}

The arithmetic operations and the elementary functions all propagate intervals:

\begin{codepairs}
\pair{arithmetic/interval-basic/1}{Interval addition: the sum of $[1,2]$ and $[3,4]$ is
$[4,6]$.}
\pair{arithmetic/interval-basic/2}{Multiplication takes the extremes over the four
corner products.}
\pair{arithmetic/interval-basic/3}{Squaring is tighter than multiplying an interval by
itself: $[-2,5]^2=[0,25]$, because the operation knows the two factors are equal.}
\pair{arithmetic/interval-basic/4}{A reciprocal that straddles zero splits into a
disjoint union of two rays---an interval may be several ranges at once.}
\pair{arithmetic/interval-basic/5}{Functions thread too: $\sqrt{[4,9]}=[2,3]$.}
\end{codepairs}

What makes these \emph{rigorous} rather than merely suggestive is \emph{outward
rounding}. Exact endpoints (integers, rationals, $\pi$) pass through untouched, but the
moment an endpoint becomes an inexact machine or MPFR real, Mathilda nudges it one unit
in the last place \emph{outward}---lower bounds down, upper bounds up---so the printed
interval is a mathematically certain enclosure, never an underestimate. There is no
dependence on an external interval library; the widening is done directly with
\texttt{nextafter} and MPFR's directed neighbours.

The classic application is certifying a floating-point computation. Recall that
$0.1+0.2-0.3$ was not exactly zero; interval arithmetic pins down exactly where the true
answer lies:

\begin{codepairs}
\pair{arithmetic/interval-rounding/1}{A certified enclosure of the rounding error: the
true value of $0.1+0.2-0.3$ lies somewhere in this tiny interval around zero.}
\pair{arithmetic/interval-rounding/2}{And zero is provably inside it---so the exact
answer really is $0$, rounding error notwithstanding.}
\end{codepairs}

A more striking use is proving that a solution \emph{exists}. If the image of a
continuous function over an interval straddles zero, the function must have a root there.
Evaluate $\sin$ over $[3, 7/2]$:

\begin{codepairs}
\pair{arithmetic/interval-root/1}{The enclosure of $\sin$ over $[3, 7/2]$, with exact
symbolic endpoints, runs from a negative value to a positive one.}
\pair{arithmetic/interval-root/2}{Zero lies inside it. This is a computer-assisted
proof that $\sin$ has a root in $[3, 7/2]$---that is, that $\pi$ lives there.}
\end{codepairs}

Interval arithmetic has one famous limitation, the \emph{dependency problem}: because
each occurrence of a variable is treated independently, an expression like $x-x$ is not
recognised as zero, and the enclosure it returns is wider than the true range. This is
not a bug but a property of the method, and it means naive interval evaluation can
over-estimate:

\begin{codepairs}
\pair{arithmetic/interval-advanced/1}{$x-x$ over $[-1,1]$ gives $[-2,2]$, not $\{0\}$:
the two copies are treated as independent.}
\pair{arithmetic/interval-advanced/2}{Iterating the logistic map $4x(1-x)$ on a point
interval, the enclosure widens at each step---a certified picture of how rounding error
grows in a chaotic system, always containing the true orbit.}
\end{codepairs}
